\chapter{{\fontsize{14pt}{21pt}\selectfont\MakeUppercase{ПРИЛОЖЕНИЕ А}\\Техническое задание}}
\label{cha:appendix-a}

\section{Общие сведения}

Настоящее техническое задание определяет требования к разработке веб-приложения на основе RAG для работы с документацией PostgreSQL с учетом версии. Техническое задание составлено для программного прототипа, реализованного в практической части выпускной квалификационной работы.

Наименование разработки: веб-приложение на основе RAG для работы с документацией PostgreSQL с учетом версии.

Основанием для разработки является задание на выпускную квалификационную работу по направлению 09.03.01 Информатика и вычислительная техника, образовательная программа <<Системная и программная инженерия>>.

Исполнитель разработки: Жатиков Артур Русланович.

\section{Назначение разработки}

Разрабатываемое веб-приложение предназначено для поиска и формирования ответов по документации PostgreSQL с учетом выбранной основной версии СУБД. Пользователь должен иметь возможность задать вопрос, выбрать версию PostgreSQL и режим ответа, после чего система должна вернуть результат с указанием источников.

Приложение ориентировано на работу с локально подготовленным корпусом документации PostgreSQL по версиям 16, 17 и 18. В кратком и подробном режимах источником фактических сведений должен быть официальный корпус документации. В обучающем режиме допускается использование дополнительного учебного корпуса как вспомогательного слоя, но официальный корпус должен оставаться основой результата.

\section{Требования к функциональным характеристикам}

Система должна поддерживать следующие пользовательские функции:

\begin{compactlist}
  \item получение списка поддерживаемых основных версий PostgreSQL;
  \item ввод пользовательского вопроса через веб-интерфейс;
  \item выбор версии PostgreSQL из поддерживаемого набора 16, 17 и 18;
  \item выбор одного из трех режимов ответа: краткого, подробного или обучающего;
  \item формирование краткого ответа по официальному корпусу документации;
  \item формирование подробного ответа по официальному корпусу документации;
  \item формирование обучающего результата со структурой краткого объяснения, предварительных условий, шагов и замечаний;
  \item отображение использованных источников с названием, адресом, типом корпуса, путем раздела, позицией в выдаче и оценкой релевантности;
  \item просмотр истории запросов и ранее сохраненных источников;
  \item отображение диагностического сообщения при ошибке обработки запроса.
\end{compactlist}

Система должна поддерживать следующие функции подготовки корпуса:

\begin{compactlist}
  \item загрузку локальных HTML-страниц официальной документации PostgreSQL по версиям 16, 17 и 18;
  \item загрузку дополнительного учебного корпуса из локальных markdown-файлов;
  \item очистку документации от служебных элементов;
  \item выделение текстовых фрагментов с сохранением заголовка, пути раздела, адреса источника, типа корпуса и версии;
  \item вычисление векторных представлений фрагментов с использованием модели \texttt{BAAI/bge-m3};
  \item сохранение документов, фрагментов и векторных представлений в PostgreSQL с расширением pgvector;
  \item административный запуск переиндексации корпуса с передачей списка версий, признаков включения корпусов и ограничения числа страниц.
\end{compactlist}

Система должна поддерживать следующие функции обработки запроса:

\begin{compactlist}
  \item валидацию вопроса, выбранной версии PostgreSQL и режима ответа;
  \item отклонение неподдерживаемой версии PostgreSQL;
  \item анализ вопроса, нормализацию текста и выделение технических терминов;
  \item построение векторного представления пользовательского вопроса;
  \item выбор допустимых корпусов по режиму ответа;
  \item поиск релевантных фрагментов в PostgreSQL с использованием pgvector;
  \item исключение официальных источников другой версии и ссылок вида \texttt{/docs/current/};
  \item повторное ранжирование найденных фрагментов по семантическим, лексическим и терминологическим признакам;
  \item проверку достаточности официальных подтверждений перед генерацией уверенного ответа;
  \item формирование безопасного отказа при недостатке подтвержденных данных;
  \item обращение к API языковой модели Groq только после подготовки проверяемого контекста;
  \item сохранение успешных и неуспешных запросов в истории.
\end{compactlist}

\section{Требования к данным и информационной совместимости}

Входными данными пользовательского запроса являются текст вопроса, выбранная основная версия PostgreSQL и режим ответа. Запрос к маршруту \texttt{/api/ask} должен содержать поля \texttt{question}, \texttt{pg\_version}, \texttt{answer\_mode}. Допустимыми значениями режима ответа являются \texttt{short}, \texttt{detailed}, \texttt{tutorial}.

Выходными данными пользовательского запроса являются режим ответа, версия PostgreSQL, текстовый ответ или обучающая структура, а также список источников. Для краткого и подробного режимов используется поле \texttt{answer}. Для обучающего режима используется объект \texttt{tutorial} с полями \texttt{short\_explanation}, \texttt{prerequisites}, \texttt{steps}, \texttt{notes}.

Данные корпуса и истории должны храниться в базе PostgreSQL. Физическая модель должна включать сущности версий, документов, фрагментов, векторных представлений, истории запросов и источников запроса. Векторные представления фрагментов должны храниться с размерностью 1024 для модели \texttt{BAAI/bge-m3}.

\section{Требования к надежности и безопасности}

Система должна предотвращать смешение официальных источников разных версий PostgreSQL. Официальный источник может использоваться в ответе только тогда, когда его адрес соответствует выбранной версии документации.

Система не должна формировать уверенный ответ, если найденный контекст не содержит достаточного официального подтверждения. В таком случае пользователь должен получить безопасную формулировку о недостатке подтвержденных данных.

Секреты внешней языковой модели и административный ключ не должны храниться в исходном коде. Они должны передаваться через переменные окружения. Контроль административного маршрута переиндексации должен выполняться при заданном административном ключе в конфигурации.

При недоступности API языковой модели система должна вернуть диагностическую ошибку и сохранить неуспешную запись истории, не подменяя ее неподтвержденным ответом.

\section{Требования к программному и техническому обеспечению}

Серверная часть должна быть реализована на Python 3.12 с использованием FastAPI, Pydantic, SQLAlchemy и Alembic. Хранение данных должно выполняться в PostgreSQL с расширением pgvector. Для построения векторных представлений должна использоваться локальная модель \texttt{BAAI/bge-m3} через библиотеку sentence-transformers. Для генерации итогового текста должен использоваться API языковой модели Groq.

Пользовательский интерфейс должен быть реализован как веб-интерфейс на HTML, CSS и JavaScript. Интерфейс должен предоставлять главную страницу для ввода запроса и страницу просмотра истории. Взаимодействие клиентской и серверной частей должно выполняться через HTTP API.

Локальное развертывание должно поддерживаться через Docker Compose. Конфигурация приложения должна задаваться переменными окружения, включая адрес базы данных, список поддерживаемых версий, параметры модели векторных представлений и ключ API языковой модели.

\section{Требования к программной документации}

В составе выпускной квалификационной работы должна быть представлена пояснительная записка, включающая анализ предметной области, постановку задачи, требования, архитектуру веб-приложения, модель данных, алгоритмы подготовки корпуса и обработки запроса, описание реализации, пользовательского интерфейса, тестирования и ограничений разработанного решения.

В приложениях должны быть представлены техническое задание, примеры запросов и ответов программного интерфейса, фрагменты листинга программного кода и дополнительные результаты тестирования.

\section{Стадии и этапы разработки}

Разработка выполняется по следующим этапам:

\begin{compactlist}
  \item анализ предметной области и структуры документации PostgreSQL;
  \item анализ проблемы учета версии PostgreSQL при формировании ответа;
  \item формулирование функциональных и нефункциональных требований;
  \item проектирование архитектуры веб-приложения и модели данных;
  \item разработка алгоритмов подготовки корпуса, поиска, повторного ранжирования и проверки подтверждений;
  \item реализация серверной части и маршрутов программного интерфейса;
  \item реализация пользовательского интерфейса и страницы истории;
  \item подготовка локального корпуса документации и дополнительного учебного корпуса;
  \item тестирование пользовательских сценариев и критичных правил обработки запроса;
  \item оформление пояснительной записки и приложений.
\end{compactlist}

\section{Порядок контроля и приемки}

Контроль выполнения требований должен проводиться на основе автоматических тестов, сценарных проверок пользовательских функций и сопоставления реализации с функциональными и нефункциональными требованиями. Критичными проверками являются валидация версии и режима ответа, выбор корпуса по режиму, исключение источников другой версии, повторное ранжирование, проверка достаточности подтверждений, сохранение истории и отображение источников.

Результат считается соответствующим техническому заданию, если приложение запускается в заданной конфигурации, принимает пользовательский запрос, учитывает выбранную версию PostgreSQL, возвращает результат выбранного режима с источниками, сохраняет историю и проходит автоматические тесты, предусмотренные практической частью.

\section{Ограничения разработки}

Разработанное приложение является инженерным прототипом. Поддерживаются только версии PostgreSQL, для которых подготовлен локальный корпус и которые указаны в конфигурации. В текущей реализации это версии 16, 17 и 18. Добавление новой версии требует наличия корпуса документации и переиндексации.

Количественные метрики качества поиска и ранжирования не нормируются в настоящем техническом задании, поскольку для них требуется отдельный размеченный набор эталонных запросов и релевантных фрагментов. Оценка результата выполняется функционально и сценарно.

\section{Трассировка требований к реализации}

Задание на выпускную квалификационную работу включено в подключаемый файл \texttt{extra/Title.pdf}. В таблице~\ref{tab:appendix-task-trace} приведена краткая трассировка ключевых пунктов задания и настоящего технического задания к фактической реализации.

\begin{center}
\small
\singlespacing
\setlength{\tabcolsep}{4pt}
\begin{longtable}{|P{0.36\textwidth}|P{0.52\textwidth}|}
\caption{Соответствие пунктов задания фактической реализации}\label{tab:appendix-task-trace}\\
\hline
\textbf{Пункт задания} & \textbf{Фактическая реализация} \\ \hline
\endfirsthead
\multicolumn{2}{l}{Продолжение таблицы \thetable}\\
\hline
\textbf{Пункт задания} & \textbf{Фактическая реализация} \\ \hline
\endhead
\endfoot
Хранение документации PostgreSQL по версиям & Таблица \texttt{versions}, локальный корпус \texttt{corpus/postgres/html/16}, \texttt{17}, \texttt{18}. \\ \hline
Разбиение документации на фрагменты & Парсер HTML и \texttt{TextChunker} с параметрами \texttt{CHUNK\_SIZE=1200}, \texttt{CHUNK\_OVERLAP=200}. \\ \hline
Векторные представления & Таблица \texttt{embeddings}, модель \texttt{BAAI/bge-m3}, размерность \texttt{1024}. \\ \hline
Поиск и повторное ранжирование & \texttt{RetrievalRepository}, \texttt{RetrievalService}, \texttt{RerankingService}. \\ \hline
Учет выбранной версии & SQL-фильтр по \texttt{Version.major\_version} и проверка URL официальных источников. \\ \hline
Режимы ответа & Схема \texttt{AskRequest.answer\_mode}, сервис \texttt{AskOrchestrationService}, интерфейс с кратким (\texttt{short}), подробным (\texttt{detailed}) и обучающим (\texttt{tutorial}) режимами. \\ \hline
Проверка подтверждений & Метод проверки достаточности контекста в \texttt{AskOrchestrationService}; безопасный отказ при недостатке официальных источников. \\ \hline
Отображение источников & Схема \texttt{SourceOut}, таблица \texttt{query\_sources}, блок источников в интерфейсе. \\ \hline
Сохранение истории & Таблицы \texttt{query\_history} и \texttt{query\_sources}, маршрут \texttt{/api/history}, страница \texttt{/history}. \\ \hline
Пользовательский интерфейс & \texttt{index.html}, \texttt{history.html}, \texttt{app.js}, \texttt{history.js}, \texttt{style.css}. \\ \hline
\end{longtable}
\end{center}

\chapter{{\fontsize{14pt}{21pt}\selectfont\MakeUppercase{ПРИЛОЖЕНИЕ Б}\\Примеры запросов и ответов программного интерфейса}}
\label{cha:appendix-b}

\section{Пример запроса в кратком режиме (\texttt{short})}

\begin{mdframed}
\small
\singlespacing
\begin{verbatim}
POST /api/ask
Content-Type: application/json

{
  "question": "Как включить логическую репликацию?",
  "pg_version": "16",
  "answer_mode": "short"
}
\end{verbatim}
\end{mdframed}

\section{Пример запроса в подробном режиме (\texttt{detailed})}

\begin{mdframed}
\small
\singlespacing
\begin{verbatim}
POST /api/ask
Content-Type: application/json

{
  "question": "Подробно объясни VACUUM в PostgreSQL",
  "pg_version": "16",
  "answer_mode": "detailed"
}
\end{verbatim}
\end{mdframed}

\section{Пример запроса в обучающем режиме (\texttt{tutorial})}

\begin{mdframed}
\small
\singlespacing
\begin{verbatim}
POST /api/ask
Content-Type: application/json

{
  "question": "Объясни EXPLAIN ANALYZE простыми словами",
  "pg_version": "16",
  "answer_mode": "tutorial"
}
\end{verbatim}
\end{mdframed}

\section{Пример ответа в кратком режиме (\texttt{short})}

\begin{mdframed}
\small
\singlespacing
\begin{verbatim}
{
  "answer_mode": "short",
  "pg_version": "16",
  "answer": "...",
  "sources": [
    {
      "title": "PostgreSQL 16 Documentation",
      "url": "https://www.postgresql.org/docs/16/...",
      "corpus_type": "official",
      "source_role": "base",
      "section_path": "Chapter 31 / Logical Replication",
      "rank_position": 1,
      "similarity_score": 0.91234
    }
  ]
}
\end{verbatim}
\end{mdframed}

\section{Пример ответа в обучающем режиме (\texttt{tutorial})}

\begin{mdframed}
\small
\singlespacing
\begin{verbatim}
{
  "answer_mode": "tutorial",
  "pg_version": "16",
  "tutorial": {
    "short_explanation": "...",
    "prerequisites": ["..."],
    "steps": ["..."],
    "notes": ["..."]
  },
  "sources": [
    {
      "title": "PostgreSQL 16 Documentation",
      "url": "https://www.postgresql.org/docs/16/...",
      "corpus_type": "official",
      "source_role": "base",
      "section_path": "Chapter 31 / Logical Replication",
      "rank_position": 1,
      "similarity_score": 0.94211
    },
    {
      "title": "Tutorial note",
      "url": "supplementary://curated/16/guide.md",
      "corpus_type": "supplementary",
      "source_role": "supplementary",
      "section_path": "Guide / Logical Replication",
      "rank_position": 2,
      "similarity_score": 0.72111
    }
  ]
}
\end{verbatim}
\end{mdframed}

\chapter{{\fontsize{14pt}{21pt}\selectfont\MakeUppercase{ПРИЛОЖЕНИЕ В}\\Руководство пользователя}}
\label{cha:appendix-c}

\section{Назначение пользовательского интерфейса}

Пользовательский интерфейс предназначен для отправки вопросов по документации PostgreSQL, выбора версии документации и выбора режима ответа. После обработки запроса приложение показывает результат, список использованных источников и сведения о последних запросах.

Главная страница приложения доступна по маршруту \texttt{/}. Страница полной истории запросов доступна по маршруту \texttt{/history}. Отдельного переключателя корпуса в интерфейсе нет: состав корпуса определяется выбранным режимом ответа.

\section{Получение ответа}

Для получения ответа пользователь выполняет следующие действия:

\begin{compactlist}
  \item открывает главную страницу приложения;
  \item вводит вопрос по документации PostgreSQL в поле запроса;
  \item выбирает поддерживаемую версию PostgreSQL: 16, 17 или 18;
  \item выбирает режим ответа: краткий, подробный или обучающий;
  \item отправляет запрос и ожидает завершения обработки;
  \item просматривает сформированный результат и список источников.
\end{compactlist}

Краткий режим предназначен для сжатого ответа по официальной документации. Подробный режим используется, когда требуется более развернутое объяснение по официальному корпусу. Обучающий режим возвращает структурированный результат: краткое объяснение, предварительные условия, пошаговые действия и замечания. В обучающем режиме дополнительный учебный корпус используется только как вспомогательный слой.

\section{Просмотр источников}

После успешной обработки запроса под ответом отображается список источников. Для каждого источника показываются название, путь раздела, тип корпуса, версия, ссылка и показатели релевантности. Если источников больше, чем помещается в начальном блоке, интерфейс позволяет раскрыть дополнительные элементы списка.

Источники нужны для ручной проверки ответа. Если система не нашла достаточного подтверждения в официальной документации выбранной версии, вместо уверенного ответа выводится безопасная формулировка о недостатке подтвержденных данных.

\section{Просмотр истории}

На главной странице отображаются последние запросы. Для полного просмотра истории пользователь переходит на страницу \texttt{/history}. На этой странице доступны вопрос, дата, выбранная версия, режим ответа, статус обработки, краткий предварительный просмотр результата и использованные источники.

История позволяет вернуться к ранее полученному ответу без повторного выполнения того же запроса. Если запрос завершился ошибкой, запись истории сохраняет статус неуспешной обработки и помогает понять, на каком этапе возникла проблема.

\section{Действия при ошибке}

Если выбранная версия не поддерживается, вопрос слишком короткий или режим ответа некорректен, приложение возвращает сообщение об ошибке валидации. Если недоступна внешняя языковая модель, пользователь получает диагностическое сообщение, а система не подменяет результат неподтвержденным ответом.

При ошибке пользователь должен проверить введенный вопрос, выбранную версию PostgreSQL и режим ответа. Если ошибка связана с внешней моделью или настройками окружения, требуется проверить конфигурацию приложения и наличие необходимых переменных окружения.

\chapter{{\fontsize{14pt}{21pt}\selectfont\MakeUppercase{ПРИЛОЖЕНИЕ Г}\\Фрагменты листинга программного кода}}
\label{cha:appendix-d}

\section{Проверка версии официального источника}

Фрагмент реализации контроля версии источника приведен в листинге~\ref{listing:appendix-version-guard}.

\captionsetup{type=listing}
\captionof{listing}{Реализация контроля версии источника}
\label{listing:appendix-version-guard}
\begin{verbatim}
@staticmethod
def _enforce_version_guard(
    *, rows: list[RetrievedChunk], pg_version: str
) -> list[RetrievedChunk]:
    expected_token = f"/docs/{pg_version}/"
    filtered: list[RetrievedChunk] = []

    for item in rows:
        if (
            item.corpus_type == CorpusType.OFFICIAL.value
            and expected_token not in item.source_url
        ):
            continue
        filtered.append(item)
    return filtered
\end{verbatim}

\section{Выбор корпусов по режиму ответа}

Фрагмент правила выбора корпусов по режиму ответа приведен в листинге~\ref{listing:appendix-corpora}.

\captionsetup{type=listing}
\captionof{listing}{Правило выбора корпусов}
\label{listing:appendix-corpora}
\begin{minted}{python}
@staticmethod
def resolve_corpora(answer_mode: str) -> list[str]:
    if answer_mode == ModeType.TUTORIAL.value:
        return [CorpusType.OFFICIAL.value, CorpusType.SUPPLEMENTARY.value]
    return [CorpusType.OFFICIAL.value]
\end{minted}

\chapter{{\fontsize{14pt}{21pt}\selectfont\MakeUppercase{ПРИЛОЖЕНИЕ Д}\\Дополнительные результаты тестирования}}
\label{cha:appendix-e}

\section{Команда запуска тестов}

\begin{mdframed}
\small
\singlespacing
\begin{minted}[breaklines=true,breakanywhere=true,fontsize=\small]{bash}
PYTHONDONTWRITEBYTECODE=1 .venv/bin/python -m pytest -q -p no:cacheprovider
\end{minted}
\end{mdframed}

\section{Результат запуска}

\begin{mdframed}
\small
\singlespacing
\begin{minted}[breaklines=true,breakanywhere=true,fontsize=\small]{text}
sss..................................................................... [ 76%]
......................                                                   [100%]
91 passed, 3 skipped in 1.00s
\end{minted}
\end{mdframed}
